随着时代的发展和项目变得越来越复杂，人们创造了一系列新的代码特性，其中一些特性在C++11和C++17中被引入，
在不熟悉这些特性的人看来，这些代码看上去显得十分无厘头，甚至有些难以理解。本节将介绍C++11和C++17中引入的一些新的代码特性，
并通过实例来说明它们的用法。

\subsubsection{lambda表达式}

\textbf{匿名函数（anonymous function）}是一种在程序中定义的函数，它没有普通函数那样的名字。C++11引入了lambda表达式，
它是一种更简洁的语法，可以用来创建匿名函数。lambda表达式的语法如下：

\begin{tcode}
常见形式：[捕获列表](参数列表) mutable noexcept -> 返回类型 {函数体}
简化版：[捕获列表](参数列表) {函数体}
\end{tcode}

\begin{enumerate}
    \item 捕获列表：捕获列表可以用来捕获外部作用域中的变量，使得lambda表达式可以访问这些变量。
    \item 参数列表：参数列表定义了lambda表达式的参数，可以有0到多个参数。
    \item mutable[可选]：mutable 允许修改值捕获的副本，不会自动修改原变量。
    \item noexcept[可选]：声明调用不会抛出异常；违反承诺会终止程序。旧的 throw(类型列表) 已在 C++17 中移除，这里不用它。
    \item 返回值：返回值定义了lambda表达式的返回类型，编译器一般可以自动推导，可以省略。
    \item 函数体：函数体定义了lambda表达式的主体，可以是一条表达式或多条语句。
\end{enumerate}

刚才的函数都起了名字，接下来会遇到“把一段行为交给另一个函数”的写法。我们把排序、捕获与移动分成独立程序，一次只看一个区别。项目放在 \path{examples/lambda-move-demo}，文件组成如下。每个 cpp 都有自己的 main，分别编译；不要把它们拼成一个文件。
\begin{lstlisting}[language={}]
lambda-move-demo/
  README.md
  capture.cpp
  forward.cpp
  function.cpp
  move-construction.cpp
  move-overloads.cpp
  mutable.cpp
  reference.cpp
  rvalue-reference.cpp
  sort.cpp
\end{lstlisting}
编译命令里，-std 指定 C++ 标准，-Wall/-Wextra 打开常用警告，-pthread 配置线程编译与链接；-o 指定临时输出文件。先进入项目目录，后面的源码名按当前示例替换。

\Needspace{9\baselineskip}
\begin{example}
用 lambda 作为比较器，对自定义结构体数组排序：
\lstinputlisting[language=C++,numbers=left,numberstyle={\color{gray}\tiny}]{examples/lambda-move-demo/sort.cpp}


从本卷目录进入 \path{examples/lambda-move-demo}，保存的源码就是 \texttt{sort.cpp}。
\begin{tbash}
g++ -std=c++17 -Wall -Wextra -pthread sort.cpp -o /tmp/cpp-demo
/tmp/cpp-demo
\end{tbash}
先输出和 3，再看到 id 依次为 7、4、2、1、-1；比较器决定了降序规则。
\end{example}

\textbf{捕获列表：}可以捕获外部作用域中的变量，包括this指针。捕获的方式包括两种，即值捕获和引用捕获。

\begin{center}
    \begin{tabular}{cccc}
        \toprule[1.5pt]
        捕获方式 & 拷贝方式 & 是否可以修改值 & 语法示例 \\
        \midrule[1pt]
        值捕获 & 保存副本 & 默认不能改副本 & [a,b]，全捕获 [=] \\
        引用捕获 & 引用原对象 & 原对象可修改时可以 & [\&a,\&b]，全捕获 [\&] \\
        \bottomrule[1.5pt]
    \end{tabular}
\end{center}

值捕获保存的是变量的副本，具体复制什么取决于类型。例如捕获一个指针，只复制地址，不会顺便复制它指向的对象。
默认的调用运算符是 const，所以不能修改捕获的普通整数副本；加上 mutable 后可以改副本，原整数仍不受影响。
引用捕获则直接使用原对象，不过你要保证调用时它还活着，不能把引用局部变量的 lambda 留到函数返回后再调用。

\Needspace{9\baselineskip}
\begin{example}
mutable 与值捕获：
\lstinputlisting[language=C++,numbers=left,numberstyle={\color{gray}\tiny}]{examples/lambda-move-demo/mutable.cpp}


从本卷目录进入 \path{examples/lambda-move-demo}，保存的源码就是 \texttt{mutable.cpp}。
\begin{tbash}
g++ -std=c++17 -Wall -Wextra -pthread mutable.cpp -o /tmp/cpp-demo
/tmp/cpp-demo
\end{tbash}
输出 1 0：被加一的是副本，外面的 a 仍是 0。
\end{example}

使用lambda表达式可以简化代码，提高可读性。由于捕获列表的存在，使得lambda表达式可以访问外部作用域的变量，
相比于普通函数，lambda表达式更加灵活。
lambda表达式的本质是一个\textbf{匿名的函数对象}（闭包），可以用 \texttt{auto} 或 \texttt{std::function} 承接；
没有任何捕获的 lambda 还可以隐式转换为函数指针，作为函数参数进行传递。

\Needspace{9\baselineskip}
\begin{example}
用 \texttt{std::function} 和 \texttt{auto} 承接 lambda：
\lstinputlisting[language=C++,numbers=left,numberstyle={\color{gray}\tiny}]{examples/lambda-move-demo/function.cpp}


从本卷目录进入 \path{examples/lambda-move-demo}，保存的源码就是 \texttt{function.cpp}。
\begin{tbash}
g++ -std=c++17 -Wall -Wextra -pthread function.cpp -o /tmp/cpp-demo
/tmp/cpp-demo
\end{tbash}
输出 3 1；两种变量都能保存可调用对象，auto 并不是“自动推导函数的返回值”。
\end{example}

在C++14和C++17中，lambda表达式增加了一些新的用法，使用这些新的特性可以简化代码，更加灵活地解决实际问题。

\textbf{泛型Lambda表达式：}
可以使用auto关键字指定形参的类型，从而快速地实现函数\textbf{泛用}。这也是\textbf{泛化编程}的一种方式，感兴趣的同学可以自行了解。
\begin{tcode}
auto f=[](auto a,auto b){return a+b;};
std::cout << f(1,2) << " " << f('a',2.0) << std::endl;
\end{tcode}

\textbf{this指针捕获：}
可以用 [this] 保存当前对象的指针，也可以从 C++17 起用 [*this] 保存当前对象的副本。两者不一样：对象后来变了，前者会看到新值，后者仍使用捕获时的副本。先看指针捕获的片段，完整对照见后面的捕获项目：
\begin{tcode}
class A{
private:
    int a = 1;
public:
    A(){
        auto function = [this](){//保存this指针，调用时对象须仍存活
            return a+1;
        };
    }
};
\end{tcode}

\textbf{初始化捕获：}在C++14中，可以使用初始化捕获来简化代码。
可以初始化一个变量，并将其绑定到lambda表达式中，这样就不需要再声明这个变量了。

\begin{tcode}
auto f = [a=1]() mutable {
    return a;
};
\end{tcode}


\Needspace{10\baselineskip}
\begin{example}
项目的 \texttt{capture.cpp} 把捕获的区别放在一起。先保存两个 lambda，再修改对象，这时候看看它们分别读到了谁：
\lstinputlisting[language=C++,numbers=left,numberstyle={\color{gray}\tiny}]{examples/lambda-move-demo/capture.cpp}
\begin{tbash}
g++ -std=c++17 -Wall -Wextra capture.cpp -o /tmp/capture-demo
/tmp/capture-demo
\end{tbash}
前两行是 2 1、2：mutable 改副本，引用捕获改原整数。最后是 5 1：保存 this 指针的 lambda 看到修改后的对象，保存 *this 副本的 lambda 看到旧值。本例对象一直活到调用结束；指针捕获不会替它延长寿命。
\end{example}
\begin{exercise}
\label{ex:v1-cpp-capture}
在 capture.cpp 中把对象的新值改为 9，先预测两种对象捕获的输出，再运行核验；解释为什么不能把保存 this 指针的 lambda 留到对象析构后使用。
\end{exercise}
\begin{exercise}\label{ex:v1-cpp-sort}

使用 \texttt{std::sort}（头文件 \texttt{<algorithm>}）配合 lambda 比较器，把一个 \texttt{std::vector<int>}
先按升序排序、再按降序排序，分别打印结果。
\end{exercise}

\subsubsection{引用和移动}
\textbf{左值与右值：}
先别急着按“等号左边、等号右边”来记，值类别描述的是\textbf{表达式}。普通变量名、解引用表达式通常是左值（lvalue），它们能指明已有对象；整数字面量和普通计算结果通常是纯右值（prvalue），例如 10 和 a+1。还有一种将亡值（xvalue），例如 \texttt{std::move(a)}：它仍指明已有对象，但允许后续操作复用其资源。纯右值与将亡值统称右值（rvalue）。

对普通整数，能否取地址可以帮助判断，但不是定义：位域是左值却不能取地址，字符串字面量也是左值；右值也不一定不可修改，例如临时类对象可以调用修改自身的成员函数。别把“临时”与“常量”混在一起。临时对象通常在完整表达式结束时销毁，满足条件的引用绑定可以延长其寿命，下面就来看这个区别。
\begin{tcode}
int a;
a = 10;//这个赋值过程中，a是左值，10是右值（字面常量）
int b;
b = a;//这个赋值过程中，b是左值，a是左值（左值可以作为右值使用）
//10 = a;//这样的语句显然是错误的，即右值不能被修改
\end{tcode}
把 a 放在赋值右边，只是在这里读取它的值，a 这个表达式本身仍是左值；值类别不由它在等号哪边决定。

\textbf{引用：}
C++的引用是一种别名，用于引用已存在的变量。引用一旦初始化，就不能再引用其他变量。引用不会为被引用对象再复制一份数据，底层是否需要保存地址由实现决定。
引用的主要用途包括函数参数传递和返回值优化。使用引用可以避免复制，提高效率，特别是在处理大型对象时。
引用通过 \& 符号来声明。例如:
\begin{tcode}
int& ref = var;
\end{tcode}
这段代码创建了一个名为 \texttt{ref} 的引用，它引用了变量 \texttt{var}。引用分为两种，左值引用（\&）和右值引用（\&\&）。
普通 \texttt{T\&} 绑定左值，\texttt{const T\&} 还可以绑定临时对象。普通 \texttt{T\&\&} 绑定右值；模板参数推导中的转发引用另有规则，稍后单独讲。

首先说左值引用，这个十分简单，可以认为是取别名，操作左值引用就相当于操作被绑定的那个左值本身，这一点就算在函数传参过程中也是如此。
\Needspace{9\baselineskip}
\begin{example}
左值引用传参：
\lstinputlisting[language=C++,numbers=left,numberstyle={\color{gray}\tiny}]{examples/lambda-move-demo/reference.cpp}


从本卷目录进入 \path{examples/lambda-move-demo}，保存的源码就是 \texttt{reference.cpp}。
\begin{tbash}
g++ -std=c++17 -Wall -Wextra -pthread reference.cpp -o /tmp/cpp-demo
/tmp/cpp-demo
\end{tbash}
输出 20，引用参数修改了 main 中的整数。
\end{example}
在这里a的值被修改为20，因为传入\texttt{func()}的a是引用，所以操作的是a本身，而不是a的拷贝。
这就是引用在传参和提高函数效率上的用处。

再说右值引用，这个就有点复杂了，右值引用是为了解决左值引用不能绑定右值的限制。
\Needspace{9\baselineskip}
\begin{example}
右值引用绑定右值：
\lstinputlisting[language=C++,numbers=left,numberstyle={\color{gray}\tiny}]{examples/lambda-move-demo/rvalue-reference.cpp}


从本卷目录进入 \path{examples/lambda-move-demo}，保存的源码就是 \texttt{rvalue-reference.cpp}。
\begin{tbash}
g++ -std=c++17 -Wall -Wextra -pthread rvalue-reference.cpp -o /tmp/cpp-demo
/tmp/cpp-demo
\end{tbash}
c 从 10 变为 20，地址由本次运行决定；这个修改本身没有发生资源移动。
\end{example}

咦？c 明明绑定了一个右值，为什么能取到地址、还能被修改？

原因在于两点：一是这里直接用临时值初始化局部右值引用，临时对象的寿命被延长到了这个引用的作用域结束；
二是\textbf{有名字的右值引用本身会被当作左值来使用}——既然有名字，就能取地址、能修改，用起来和普通左值引用没有区别。

那右值引用存在的意义是什么？设想一个超级大的对象：传出函数时，我们希望把它的资源交给新对象，避免不必要的复制。实际返回还可能经过返回值优化，不能把所有返回都算成复制。
右值引用让我们可以直接"接管"临时对象的内存，避免这次拷贝——具体能否转移、怎样转移，由类型的移动构造或移动赋值实现决定。仅仅绑定右值引用，并不会替你搬数据。
这就是移动构造函数存在的理由，马上就能在下面的例子里看到。

\Needspace{9\baselineskip}
\begin{example}
一个演示默认、含参、拷贝、移动构造与拷贝赋值的类（不代表写全了所有特殊成员函数），请仔细体会每一行的输出：
\lstinputlisting[language=C++,numbers=left,numberstyle={\color{gray}\tiny}]{examples/lambda-move-demo/move-construction.cpp}


从本卷目录进入 \path{examples/lambda-move-demo}，保存的源码就是 \texttt{move-construction.cpp}。
\begin{tbash}
g++ -std=c++17 -Wall -Wextra -pthread move-construction.cpp -o /tmp/cpp-demo
/tmp/cpp-demo
\end{tbash}
对照上面的执行顺序检查构造、赋值和析构记录。
\end{example}
这时候回头看 a5 和 a6：a5 只是给 a1 多取了一个别名，没有调用移动构造；a6 才创建新对象并调用移动构造。移动后的字符串仍有效，但内容不保证一定为空，不要把某次输出当成固定规则。

\textbf{移动语义：}
首先了解一个std中的函数move，它的作用是将参数类型强制转换为右值，不管参数原来是左值还是右值。
\Needspace{9\baselineskip}
\begin{example}
\texttt{std::move} 的效果：
\lstinputlisting[language=C++,numbers=left,numberstyle={\color{gray}\tiny}]{examples/lambda-move-demo/move-overloads.cpp}


从本卷目录进入 \path{examples/lambda-move-demo}，保存的源码就是 \texttt{move-overloads.cpp}。
\begin{tbash}
g++ -std=c++17 -Wall -Wextra -pthread move-overloads.cpp -o /tmp/cpp-demo
/tmp/cpp-demo
\end{tbash}
调用类别依次是右值、左值、右值、const 左值、const 右值。
\end{example}
如代码所见，当我们使用std::move()时，编译器将左值转换为右值。

再介绍一种利用模板实现完美转发的方式。在模板中，T\&\&可以接收左值和右值，要利用到forward<>()函数，它会保留引用的类型，并将其转发给另一个函数，
因为有名字的参数 t 在函数体中是左值表达式，直接传 t 会选择左值重载。这里的 T 必须是被推导的模板参数；并非所有 T\&\& 都是转发引用。

\Needspace{9\baselineskip}
\begin{example}
完美转发：
\lstinputlisting[language=C++,numbers=left,numberstyle={\color{gray}\tiny}]{examples/lambda-move-demo/forward.cpp}


从本卷目录进入 \path{examples/lambda-move-demo}，保存的源码就是 \texttt{forward.cpp}。
\begin{tbash}
g++ -std=c++17 -Wall -Wextra -pthread forward.cpp -o /tmp/cpp-demo
/tmp/cpp-demo
\end{tbash}
先运行原版，再按紧邻的提示替换 Fun(t)，比较两次记录。只有后者会保留右值调用类别。
\end{example}

\begin{note}
先运行 \texttt{Fun(t)}，再把这一行换成 \texttt{Fun(std::forward<T>(t))}，不要两行同时保留。对比调用记录：前者把有名字的参数作为左值传入，后者保留传入时的值类别。引用折叠是类型规则，与表达式有名字后成为左值要分开看。
\end{note}

\begin{conclusion}
左值、右值与两种引用的对照：

\begin{itemize}
\item 值类别属于表达式：变量名通常是左值，10 是纯右值，std::move(x) 是将亡值；不要只按取地址或等号位置判断；
\item 普通 T\& 绑定左值，const T\& 还能绑定临时对象；引用必须在对象有效期内使用；
\item 普通 T\&\& 绑定右值，有名字的右值引用作为表达式仍是左值；绑定本身不转移资源；
\item \texttt{std::move}：把左值强制转成右值（本身不移动任何数据），\texttt{std::forward}：在模板中保留参数的左右值属性。
\end{itemize}
\end{conclusion}

左值引用和右值引用的概念比较容易混淆，要在实践中不断运用，才能熟练掌握。

\begin{exercise}\label{ex:v1-cpp-reference}

编写一个函数 \texttt{void swapVal(int\& a, int\& b)} 交换两个整数的值，在 \texttt{main} 中调用并验证；
再把参数类型改为按值传递（去掉 \texttt{\&}），观察交换是否仍然生效，解释原因。再对照 move-construction.cpp 的 a5、a6，说明哪一处真正调用了移动构造；按提示切换 forward.cpp 的调用语句，预测并核验五次重载选择。
\end{exercise}
